% Ad Familiares 10.28 --- headnote
% ad-familiares-10-28, ca. 2 February 43 BC, Rome

\textit{Cicero to C. Trebonius, written from Rome about 2
February 43 BC --- Perseus dateline \textit{Scr. Romae circ.
iv Non. Febr. a. 711 (43)}. Trebonius, one of the conspirators
of 15 March 44, is now governor of the province of Asia; he
will be killed by Dolabella at Smyrna within weeks of this
letter's arrival. The opening sentence is the most famous in
the corpus of letters of this period and one of the most cited
in all of Cicero: \textit{quam vellem ad illas pulcherrimas
epulas me Idibus Martiis invitasses! reliquiarum nihil
haberemus} --- ``How I wish you had invited me to that most
beautiful banquet on the Ides of March! There would be no
leftovers.'' The metaphor is sustained: the banquet was the
killing of Caesar, the leftovers are Antony, and Cicero's
retrospective regret is that the conspirators, by sparing the
consul, left him the political work of finishing what they
had begun. Trebonius is gently and unmistakably reproached
in the next sentence for personally drawing Antony aside on
the day so that he should not enter the senate house ---
\textit{quod vero a te, viro optimo, seductus est tuoque
beneficio adhuc vivit haec pestis}.}

\textit{The body of the letter is a brisk political report
written for a man who has been months in his province. Cicero
dates his own resumption of leadership precisely to 20
December 44 --- the meeting of the Senate, called by the
tribunes, at which he delivered the Third Philippic and
launched the senatorial war against Antony. ``That day and my
exertion and action first brought the Roman people the hope
of recovering its liberty'' is Cicero's own settled account
of where the resistance began, and he gives it to Trebonius
in the same week that the Senate has finally declared Antony
a public enemy. The roll call at section 3 is characteristic
end-of-letter shorthand: a brave Senate, mixed consulars,
the loss of Servius Sulpicius Rufus (who died on his embassy
to Antony in late January and whose funeral oration is the
Ninth Philippic), L. Iulius Caesar restrained by being
Antony's uncle, ``the consuls excellent, D. Brutus splendid,
the boy Caesar admirable.'' The qualification on Octavian ---
\textit{de quo spero equidem reliqua} --- is the formal hope
that he will continue as he has begun; the candid note that
follows, that without his hasty enlistment of the veterans
``no crime, no cruelty would Antony have left undone,'' shows
how unreservedly Cicero is, in early February, still backing
the young Caesar against Antony. By the time Trebonius could
have read this, both Trebonius and the diagnosis of
section 3 were obsolete.}
